La lesió per sesamoïditis (inflamació de l'os sesamoide del peu) és relativament habitual entre la població que
practica esports d'impacte (atletisme, bàsquet, tennis\ldots). En una població d'esportistes, s'ha fet un estudi
diferenciant entre els que practiquen esports d'impacte i els que practiquen esports sense impacte brusc (com
ara natació, pilates, senderisme\ldots). S'ha pogut determinar que el 45\,\% practiquen esports d'impacte. Entre
aquests, un 10\,\% pateixen lesions per sesamoïditis, mentre que entre els que no practiquen esports d'impacte
només un 3\,\% presenten aquesta lesió. Escollim un esportista a l'atzar.

\begin{apartats}

\apartat{0,75}
Quina és la probabilitat que pateixi sesamoïditis?

\begin{solucio}
Considerem els esdeveniments aleatoris $I$ = «practicar un esport d'impacte» i $S$ = «patir sesamoïditis». De
les dades de l'enunciat sabem que $P(I)=0{,}45$, $P\bigl(\overline{I}\bigr)=1-0{,}45=0{,}55$, $P(S\mid I)=0{,}1$ i
$P\bigl(S\mid\overline{I}\bigr)=0{,}03$. Per la llei de les probabilitats totals, tenim
\[
P(S)=P(S\mid I)\,P(I)+P\bigl(S\mid\overline{I}\bigr)\,P\bigl(\overline{I}\bigr)=0{,}1\cdot0{,}45+0{,}03\cdot0{,}55=0{,}0615 .
\]

\textit{Pauta oficial:} 0,25 per plantejar la llei de la probabilitat total (o fer un arbre de decisió o
similar); 0,25 per identificar correctament les dades del problema, i 0,25 pel càlcul.
\end{solucio}

\apartat{0,75}
Si l'esportista escollit té una lesió per sesamoïditis, quina és la probabilitat que practiqui esports
d'impacte?

\begin{solucio}
Per la fórmula de Bayes,
\[
P(I\mid S)=\frac{P(S\mid I)\,P(I)}{P(S)}=\frac{0{,}1\cdot0{,}45}{0{,}0615}\approx0{,}7317\ldots
\]

\textit{Pauta oficial:} 0,25 per plantejar correctament la fórmula de Bayes (o la fórmula de la probabilitat
condicionada) i 0,5 pel càlcul.
\end{solucio}

\apartat{1}
Una empresa de calçat esportiu ha creat una sabatilla amb amortiment per a minimitzar les lesions per
sesamoïditis. Els beneficis generats per la venda d'aquest producte, en milers d'euros, segueixen una funció
de la forma $f(x)=ax^3+bx^2+cx$, on $x$ són els anys transcorreguts des que la sabatilla va sortir a la venda i
$a$, $b$ i $c$ són constants reals.

Calculeu els valors de $a$, $b$ i $c$ sabent que el primer any es va obtenir el màxim de beneficis, amb un
valor de 8\,000 euros, i que el segon any va haver-hi un punt d'inflexió en els beneficis.

\begin{solucio}
El primer any s'obté el màxim de beneficis; per tant, $x=1$ és un màxim de la funció i $f'(1)=0$. Aquests
beneficis són de 8\,000 €; per tant, $f(1)=8$. També ens diuen que el segon any hi va haver un punt d'inflexió
en els beneficis; per tant, $f''(2)=0$. Calculem la primera i la segona derivades de $f(x)=ax^3+bx^2+cx$,
$f'(x)=3ax^2+2bx+c$ i $f''(x)=6ax+2b$, i imposem aquestes tres condicions:
\[
\left.\begin{aligned}f'(1)=0&\Rightarrow3a+2b+c=0\\f(1)=8&\Rightarrow a+b+c=8\\f''(2)=0&\Rightarrow12a+2b=0\end{aligned}\right\}
\]
Restant les dues primeres equacions, obtenim $2a+b=-8$ i, combinant-ho amb la tercera,
\[
\left.\begin{aligned}12a+2b&=0\\2a+b&=-8\end{aligned}\right\}\;\Rightarrow\;
\left.\begin{aligned}12a+2b&=0\\-4a-2b&=16\end{aligned}\right\}\;\Rightarrow\;8a=16\;\Rightarrow\;a=\frac{16}{8}=2 .
\]
Per tant, $2\cdot2+b=-8$, d'on $b=-12$. Finalment, $a+b+c=8\Rightarrow c=8-a-b=8-2+12=18$. La funció que
compta els beneficis d'aquesta empresa és, doncs, $f(x)=2x^3-12x^2+18x$.

\textit{Pauta oficial:} 0,25 per imposar $f'(1)=0$, 0,25 per imposar $f(1)=8$, 0,25 per imposar $f''(2)=0$ i
0,25 pel càlcul final.
\end{solucio}

\end{apartats}
